
\section{Annexos}

\subsection{Codi Font}

\subsubsection{main}

\begin{verbatim}

% PIPELINE COMPLETO DE DETECCIÓN Y RECONOCIMIENTO DE MATRÍCULAS

addpath(genpath(pwd));

function image = cutImage(im, ROI)
    image = imcrop(im, ROI);
end


% Flags para controlar la visualización y ejecución del sistema

showSteps = false;      % Muestra pasos intermedios del detector de ROIs
showSegments = true;        % Muestra segmentación de caracteres dentro de la ROI

loadDataset = false;
useTxtFiles = true;
fileFactor = 1;
showVisuals = true;

% Ruta del dataset
datasetPath = fullfile(pwd, 'Datasets', 'Dataset');

% Lista de imagenes a procesar
images = ["test_092.jpg","test_071.jpg","test_073.jpg","test_061.jpg","test_057.jpg","test_060.jpg","test_048.jpg","test_046.jpg","test_023.jpg","test_017.jpg","test_015.jpg","test_010.jpg","eu8.jpg","eu4.jpg","eu11.jpg",];

% Carga del dataset
if useTxtFiles
    disp('Using txt');
    files = dir(fullfile(datasetPath, '*.txt'));
    images = [];
    roisDataset = [];
    platesDataset = [];
    for f = 1:(numel(files) * fileFactor)
        file = fullfile(datasetPath, files(f).name);
        fid = fopen(file, 'r');
        data = textscan(fid, '%s %d %d %d %d %s');
        fclose(fid);

        if ~isempty(data{1})
            imgPath = fullfile(datasetPath, string(data(1)));
            x = data{2};
            y = data{3};
            w = data{4};
            h = data{5};
            pl = string(data{6});
            bb = [x, y, w, h];
            images = [images(:); imgPath(:)];
            roisDataset = [roisDataset; bb];
            platesDataset = [platesDataset; pl];
        end
    end
elseif loadDataset
    images = [];
    data = [
        dir(fullfile(datasetPath, '*.jpg'));
        dir(fullfile(datasetPath, '*.png'));
        dir(fullfile(datasetPath, '*.jpeg'))
    ];

    images = fullfile({data.folder}, {data.name});
    images = string(images);

else 
    images = fullfile(datasetPath, images);
    images = string(images);
end


% Bucle principal de procesamiento
exactas = 0;
parciales = 0;
errores = 0;

confusionLog = {};

for k = 1:numel(images)

    % Lectura de la imagen actual
    im = imread(images(k));
    if useTxtFiles
        expected = char(platesDataset(k));
        bestMatch = 0;
        bestPlate = "";
    end

    % FASE 1: DETECCIÓN DE POSIBLES MATRÍCULAS (ROIs)
    candidates = getCandidates(im, false);
    
    % Clasificar todas las ROIs y recoger scores ANTES del NMS
    scores = zeros(size(candidates, 1), 1);
    validMask = false(size(candidates, 1), 1);
    
    for i = 1:size(candidates, 1)
        crop = cutImage(im, candidates(i,:));
        [isPlate, scores(i)] = classify_roi(crop);
        validMask(i) = isPlate;
    end
    
    % Quedarse solo con las que el clasificador acepta
    candidates = candidates(validMask, :);
    scores = scores(validMask);
    
    % NMS sobre las candidatas válidas
    candidates = nms(candidates, scores, 0.45, 0.75);
    
    [~, name] = fileparts(images(k));
    if showVisuals
        og = figure('Name', name);
        imshow(im), title(name);
        hold on;
    end
    
    % FASE 2: ALINEACIÓN Y SEGMENTACIÓN
    for i = 1:size(candidates, 1)
        plate = cutImage(im, candidates(i,:));
        match = 0;
    
        [aligned, BB] = alignPlate(plate, candidates(i,:));
        aligned = imresize(aligned, 2);
        blobs = segm(aligned);

        [plateText, charLabels] = recognize_plate(aligned, blobs);
        %fprintf("Matrícula detectada: %s\n", plateText);

        if size(blobs,1) <= 3
            [aligned, BB] = alignPlate(255-plate, candidates(i,:));
            aligned = imresize(aligned, 2);
            blobs = segm(aligned);
            [plateText, charLabels] = recognize_plate(aligned, blobs);
            %fprintf("Matrícula detectada: %s\n", plateText);
        end
        %fprintf("Matrícula detectada: %s\n", plateText);
        detectedCh = char(plateText);

        if useTxtFiles
            if strcmp(detectedCh,expected)
                match = 2;
            else
                m = 0;
                tmp = expected;
                for c = 1:length(detectedCh)
                    idx = strfind(tmp, detectedCh(c));
                    if ~isempty(idx)
                        m = m + 1;
                        tmp(idx(1)) = [];
                    end
                end
                if m >= 3
                    match = 1;
                end
            end
    
            if match > bestMatch
                bestMatch = match;
                bestPlate = plateText;
            end
        end
    
        % Visualización

        if showVisuals
            figure(og);
            if useTxtFiles
                bbDataset = roisDataset(k, :);
                rectangle('Position',bbDataset, 'EdgeColor', 'g', 'LineWidth', 2.5);
            end

            if ~isempty(BB)
                x_pl = [BB(:,1); BB(1,1)];
                y_pl = [BB(:,2); BB(1,2)];
                line(x_pl, y_pl, 'Color', 'r', 'LineWidth', 2.5);
        
                if size(blobs,1) > 3 && showSegments
                    figure, imshow(aligned), title('cropped image');
                    hold on;
                    for j = 1:size(blobs,1)
                        bb = blobs(j,:);
                        rectangle('Position', bb, ...
                            'EdgeColor', 'r', ...
                            'LineWidth', 2);
                        text(bb(1)+bb(3)/2, ...
                             bb(2)-5, ...
                             char(charLabels(j)), ...
                             'FontSize', 16, ...
                             'FontWeight', 'bold', ...
                             'Color', 'r', ...
                             'HorizontalAlignment', 'center');
                    end
                    if showVisuals
                        hold off;
                    end
                end
            end
        end
    end

    if showVisuals
        figure(og);
        hold off;
    end

    if useTxtFiles
        if bestMatch == 2
            exactas = exactas + 1;
            ver = "Correcto";
        elseif bestMatch == 1
            parciales = parciales + 1;
            ver = "Parcialmente correcto";

            % Analisis de Confusiones
            det = char(bestPlate);
            exp = expected;
            tmpExp = exp;
            tmpDet = det;
            % Elimina los caracteres que coinciden para quedarse solo con errores
            for c = length(tmpDet):-1:1
                idx = strfind(tmpExp, tmpDet(c));
                if ~isempty(idx)
                    tmpExp(idx(1)) = [];
                    tmpDet(c) = [];
                end
            end
            % Empareja lo que faltaba con lo que se detectó de más
            nPairs = min(length(tmpExp), length(tmpDet));
            for c = 1:nPairs
                confusionLog{end+1} = {tmpExp(c), tmpDet(c)};
            end
        else
            errores = errores + 1;
            ver = "Error";
        end
        fprintf("Archivo: %s\n", name);
        fprintf("Matrícula detectada: %s\n", bestPlate);
        fprintf("Matrícula esperada:  %s\n", expected);
        fprintf("Veredicto: %s\n", ver);
        fprintf("=======================\n");
    end
    if ~showVisuals
        close all hidden;
    end
end

fprintf("\n===============================\n");
fprintf("     RESULTADOS FINALES      \n");
fprintf("Imagenes evaluadas: %d\n", numel(images));
fprintf("===============================\n");
fprintf("Iguales: %d\n", exactas);
fprintf("Parciales: %d\n", parciales);
fprintf("Errores: %d\n", errores);
fprintf("===============================\n");
precision = (exactas / numel(images));
precParcial = ((exactas + parciales) / numel(images));
fprintf("Tasa de acierto: %.0f%%\n", precision * 100);
fprintf("Tasa de acierto parcial: %.0f%%\n",  precParcial * 100);

% RESUMEN DE CONFUSIONES 
if ~isempty(confusionLog)
    % Agrupa y cuenta cada par (real -> detectado)
    pares = cellfun(@(x) sprintf('%s->%s', x{1}, x{2}), confusionLog, 'UniformOutput', false);
    [paresSorted, ~, ic] = unique(pares);
    counts = accumarray(ic, 1);
    [countsSorted, sortIdx] = sort(counts, 'descend');

    fprintf("\n=== CONFUSIONES MÁS FRECUENTES ===\n");
    fprintf("  Real -> Detectado   Veces\n");
    fprintf("  --------------------------\n");
    for i = 1:length(paresSorted)
        fprintf("  %-18s  %d\n", paresSorted{sortIdx(i)}, countsSorted(i));
    end
    fprintf("===================================\n");
end
    
\end{verbatim}

\subsubsection{candidates\_detection}

\begin{verbatim}

function candidates = getCandidates(im, showSteps)
    im = rgb2gray(im);
    im = imadjust(im);
    [szy, szx] = size(im);
    if showSteps
        figure, imshow(im), title('input image');
    end
    candidates = [];
    
    area = numel(im);
    edges = edge(im,'sobel', 'vertical');
    if showSteps
        figure, imshow(edges), title('edges');
    end
    sizes = [
        1, 5;
        2, 12;
        5, 20;
        8, 30;
        12, 45;
        ceil(0.010 * szy), ceil(0.030 * szx);
        ];
    for si = 1:size(sizes, 1)
        he = sizes(si, 1);
        wi = sizes(si, 2);
        EE = strel('rectangle', [he, wi]);
        EE2 = strel('rectangle', [he, wi]);
    
        cl = imdilate(edges, EE);
        cl = imfill(cl, 'holes');
        cl = imerode(cl, EE2);
    
        if showSteps
            figure, imshow(cl), title('close');
        end
    
        CC = bwconncomp(cl);
        S = regionprops(CC, 'BoundingBox', 'Area', 'Extent'); 

        for k = 1:numel(S)
            bb = S(k).BoundingBox;
            x = max(bb(1) - 10, 0);
            y = max(bb(2) - 10, 0);
            w = min(bb(3) + 20, szx);
            h = min(bb(4) + 20, szy);
            if w < 10 || h < 10
                continue;
            end
    
            ar = w / h;
            if ar < 1
                ar = 1 / ar;
            end
    
            areaRatio = S(k).Area / area;
    
            isWide = (ar > 1.7) && (ar <= 12);
            isRect = S(k).Extent >= 0.15;
            isReasonable = (areaRatio > 0.0005) & (areaRatio <= 0.25);
            if isWide && isRect && isReasonable
                candidates = [candidates; [x y w h]];
            end
        end
    end
end

\end{verbatim}

\begin{verbatim}

function kept = nms(candidates, scores, iouThresh, containThresh)

    if nargin < 4, containThresh = 0.75; end
    if nargin < 3, iouThresh = 0.45; end

    if isempty(candidates)
        kept = zeros(0, 4);
        return;
    end

    boxes  = double(candidates);
    scores = double(scores);

    % Ordenar por score descendente
    [~, order] = sort(scores, 'descend');
    boxes = boxes(order, :);

    keep = true(size(boxes, 1), 1);

    for i = 1:size(boxes,1)-1
        if ~keep(i)
            continue; 
        end
        for j = i+1:size(boxes,1)
            if ~keep(j)
                continue; 
            end

            [iou, r1, r2] = getOverlap(boxes(i,:), boxes(j,:));

            if iou >= iouThresh || r1 >= containThresh || r2 >= containThresh
                keep(j) = false;  % eliminar la de menor score
            end
        end
    end
    kept = boxes(keep, :);
end


function [iou, r1, r2] = getOverlap(A, B)
    x1 = max(A(1), B(1));
    y1 = max(A(2), B(2));
    x2 = min(A(1)+A(3), B(1)+B(3));
    y2 = min(A(2)+A(4), B(2)+B(4));

    iw = max(0, x2 - x1);
    ih = max(0, y2 - y1);
    intersection = iw * ih;

    aA = A(3) * A(4);
    aB = B(3) * B(4);
    union = aA + aB - intersection;

    iou = intersection / union; 
    r1  = intersection / aA;      % Fracción de A cubierta por B
    r2  = intersection / aB;      % Fracción de B cubierta por A
end

\end{verbatim}

\subsubsection{plate\_classifier}

\begin{verbatim}

function [prediction, score] = classify_roi(img)

    % Load model
    
    persistent model;
    
    if isempty(model)
    
        data = load('plate_classifier.mat');
    
        model = data.model;
    
    end
    
    % Config
    
    imageSize = [64 128];
    
    % Features
    
    features = extract_plate_features(img, imageSize);
    
    % Prediction
    
    [prediction, scores] = predict(model, features);
    
    % Score
    
    score = scores(2);

end
    
\end{verbatim}


\begin{verbatim}

function features = extract_plate_features(img, imageSize)
    
    % Grayscale
    
    if size(img,3) == 3
        img = rgb2gray(img);
    end

    % Aspect and Area Ratio

    [h, w] = size(img);
    aspectRatio = w / h;
    areaRatio = (h * w);
    
    % Resize
    
    img = imresize(img, imageSize);
    img = im2uint8(im2double(img));
    
    % HOG FEATURES
    
    hog_feature = extractHOGFeatures(img,'CellSize', [8 8]);

    features = [hog_feature, aspectRatio, areaRatio] ;

end
    
\end{verbatim}

\begin{verbatim}

clear;
clc;
close all;

%% Configuration

positiveDir = 'dataset/positive';
negativeDir = 'dataset/negative';

imageSize = [64 128]; % height * width

%% Load dataset

positiveImages = dir(fullfile(positiveDir, '*.jpg'));
negativeImages = dir(fullfile(negativeDir, '*.jpg'));

numPos = length(positiveImages);
numNeg = length(negativeImages);

fprintf('Positives: %d\n', numPos);
fprintf('Negatives: %d\n', numNeg);

%% Features extraction

features = [];
labels = [];


% Positives

for i = 1:numPos

    imgPath = fullfile(positiveDir, positiveImages(i).name);

    img = imread(imgPath);

    feat = extract_plate_features(img, imageSize);

    features = [features; feat];

    labels = [labels; 1];

end


% Negatives

for i = 1:numNeg

    imgPath = fullfile(negativeDir, negativeImages(i).name);

    img = imread(imgPath);

    feat = extract_plate_features(img, imageSize);

    features = [features; feat];

    labels = [labels; 0];

end

fprintf('Features extracted\n');
size(features)

%% Train / Test division

rng(42);
cv = cvpartition(labels, 'HoldOut', 0.2);

Xtrain = features(training(cv), :);
Ytrain = labels(training(cv));

Xtest = features(test(cv), :);
Ytest = labels(test(cv));

%% Train classifier (SVM)

fprintf('Training SVM...\n');

model = fitcsvm(Xtrain, Ytrain, 'KernelFunction', 'linear', ...
    'Standardize', true, 'ClassNames', [0 1], 'KernelScale','auto');

fprintf('Model trained\n');

%% Evaluation

predictions = predict(model, Xtest);

accuracy = sum(predictions == Ytest) / length(Ytest);

fprintf('Accuracy: %.2f %%\n', accuracy * 100);

%% Confusion matrix

figure;
confusionchart(Ytest, predictions);

title('Confusion Matrix');

%% Save model

save('plate_classifier.mat', 'model');

fprintf('Model saved in plate_classifier_model.mat\n');
    
\end{verbatim}


\subsubsection{character\_segmentation}

\begin{verbatim}


function [alignedCrop, alignedBB] = alignPlate(im, boundingBox)
    if size(im, 3) == 3
        im = rgb2gray(im);
    end
    bw = imadjust(im);
    bw = imbinarize(bw, "adaptive");
    [szy, szx] = size(im);
    alignedBB = boundingBox;
    
    CC = bwconncomp(bw);
    S = regionprops(CC, 'Orientation', 'Area', 'Extent');
    [~, idx] = max([S.Area]);
    if ~isempty(S)
        ang = S(idx).Orientation;
        im = imrotate(im, -ang);
        bw = imrotate(bw, -ang);
    
        CC = bwconncomp(bw);
        S = regionprops(CC, 'BoundingBox', 'Area');
        [~, idx] = max([S.Area]);
        bb = S(idx).BoundingBox;
    
        [szyR, szxR] = size(im);
        
        off = 15;
        bx = max(1, bb(1) - off);
        by = max(1, bb(2) - off);
        bwi = min(szxR - bx, bb(3) + (2 * off));
        bhe = min(szyR - by, bb(4) + (2 * off));
        bb2 = [bx, by, bwi, bhe];

        im = imcrop(im, bb);

        rot = [
            bb(1), bb(2);
            bb(1) + bb(3), bb(2);
            bb(1) + bb(3), bb(2) + bb(4);
            bb(1), bb(2) + bb(4);
            ];
    
    
        CO = [szx, szy] / 2;
        CR = [szxR, szyR] / 2;
        R = [cosd(ang), -sind(ang); sind(ang), cosd(ang)];
        corners = zeros(4, 2);
        for i = 1:4
            pR = rot(i, :) - CR;
            pO = pR * R;
            corners(i, :) = pO + CO;
        end
    
        alignedBB = corners + [boundingBox(1), boundingBox(2)] - 1;
    
    end
    alignedCrop = im;
    %figure, imshow(alignedCrop);
end

\end{verbatim}

\begin{verbatim}


function chars = segm(im)
   
    if size(im, 3) == 3
        im = rgb2gray(im);
    end

    im = imgaussfilt(im, 0.5);
    bw = ~imbinarize(im, 'adaptive', 'ForegroundPolarity', 'dark', 'Sensitivity', 0.5);

    bw = imclearborder(bw); % Elimina objetos tocando los bordes
    bw = bwareaopen(bw, 30); 

    CC = bwconncomp(bw);

    % Close para conectar numeros como 3, etc
    se = strel('square', 2); 
    bw = imclose(bw, se);

    S = regionprops(CC, 'BoundingBox', 'Area', 'Extent');
    
    [imgHeight, ~] = size(bw);
    blobs = [];

    for k = 1:numel(S)
        bb = S(k).BoundingBox;
        x = bb(1); 
        y = bb(2); 
        w = bb(3); 
        h = bb(4);

        ar = w / h; % Aspect Ratio

        cond_ar = (ar > 0.05 && ar < 1.5);

        cond_h = (h > imgHeight * 0.35 && h < imgHeight * 0.95);

        % Extent (Area del objeto / Area del Bounding Box)
        cond_ext = (S(k).Extent > 0.12 && S(k).Extent < 1.0);

        if cond_ar && cond_h && cond_ext
            dfW = w * 0.10;
            dfH = h * 0.05;
            x2 = x - dfW;
            y2 = y - dfH;
            w2 = w + dfW * 2;
            h2 = h + dfH * 2;
            blobs = [blobs; [x2, y2, w2, h2]];
        end
    end

    % Para descartar simbolos/banderas
    if size(blobs, 1) >= 3
        medH = median(blobs(:, 4));
        keepMask = blobs(:, 4) > medH * 0.80;
        blobs = blobs(keepMask, :);
    end

    if ~isempty(blobs)
        chars = sortrows(blobs, 1);
    else
        chars = [];
    end
end

\end{verbatim}


\subsubsection{dataset\_generation}

\begin{verbatim}

% Intersection over union  
% Checks how overlapped is the negative crop compared with the license
% plate

function iou = compute_iou(boxA, boxB)

    xA = max(boxA(1), boxB(1));
    yA = max(boxA(2), boxB(2));
    
    xB = min(boxA(1)+boxA(3), boxB(1)+boxB(3));
    yB = min(boxA(2)+boxA(4), boxB(2)+boxB(4));
    
    interW = max(0, xB - xA);
    interH = max(0, yB - yA);
    
    interArea = interW * interH;
    
    areaA = boxA(3) * boxA(4);
    areaB = boxB(3) * boxB(4);
    
    unionArea = areaA + areaB - interArea;
    
    iou = interArea / unionArea;

end
    
\end{verbatim}

\begin{verbatim}

clear;
clc;
close all;

%% Configuration

imagesDir = 'images';
annotationsDir = 'annotations';

positiveDir = 'dataset/positive';
negativeDir = 'dataset/negative';

mkdir(positiveDir);
mkdir(negativeDir);

padding = 0.05; % 5%

negativePerImage = 5;

%% Read Images

imageFiles = dir(fullfile(imagesDir, '*.png'));

positiveCount = 1;
negativeCount = 1;

%% Processing

for i = 1:length(imageFiles)

    fprintf('Processing image %d / %d\n', i, length(imageFiles));

    % Read Image

    imageName = imageFiles(i).name;

    imgPath = fullfile(imagesDir, imageName);

    img = imread(imgPath);

    [imgH, imgW, ~] = size(img);

    % Read XML

    xmlName = replace(imageName, '.png', '.xml');

    xmlPath = fullfile(annotationsDir, xmlName);

    xmlDoc = xmlread(xmlPath);

    objects = xmlDoc.getElementsByTagName('object');

    if objects.getLength == 0
        continue;
    end

    % Read only license plate

    obj = objects.item(0);

    bbox = obj.getElementsByTagName('bndbox').item(0);

    xmin = str2double(bbox.getElementsByTagName('xmin').item(0).getTextContent);

    ymin = str2double(bbox.getElementsByTagName('ymin').item(0).getTextContent);

    xmax = str2double(bbox.getElementsByTagName('xmax').item(0).getTextContent);

    ymax = str2double(bbox.getElementsByTagName('ymax').item(0).getTextContent);

    % Add padding

    bw = xmax - xmin;
    bh = ymax - ymin;

    padX = round(bw * padding);
    padY = round(bh * padding);

    xminP = max(1, xmin - padX);
    yminP = max(1, ymin - padY);

    xmaxP = min(imgW, xmax + padX);
    ymaxP = min(imgH, ymax + padY);

    % Positive crop

    positiveCrop = img(yminP:ymaxP, xminP:xmaxP, :);

    positiveName = sprintf('pos_%05d.jpg', positiveCount);

    positiveCrop = im2uint8(mat2gray(positiveCrop));

    imwrite(positiveCrop, fullfile(positiveDir, positiveName));

    positiveCount = positiveCount + 1;

    % Generate negatives (no plate)

    for n = 1:negativePerImage

        validCrop = false;

        attempts = 0;

        while ~validCrop && attempts < 100

            attempts = attempts + 1;

            % Similar size to the license plate

            cropW = bw;
            cropH = bh;

            rx = randi([1, imgW - cropW]);
            ry = randi([1, imgH - cropH]);

            % IOU

            iou = compute_iou([rx ry cropW cropH], [xmin ymin bw bh]);

            if iou < 0.1

                negativeCrop = img(ry:ry+cropH, rx:rx+cropW, :);

                negativeName = sprintf('neg_%05d.jpg', negativeCount);

                negativeCrop = im2uint8(mat2gray(negativeCrop));

                imwrite(negativeCrop, fullfile(negativeDir, negativeName));

                negativeCount = negativeCount + 1;

                validCrop = true;

            end

        end

    end

end

fprintf('Dataset successfully generated\n');
    
\end{verbatim}


\subsubsection{ocr\_classifier}

\begin{verbatim}

function label = classify_character(img)

    % Load model
    
    persistent svmModel5;
    
    if isempty(svmModel5)
    
        data = load('ocr_model5.mat');
    
        svmModel5 = data.svmModel5;
    
    end

    % Config 

    imgSize = [32 32];
    
    % Extract Features

    features = extract_character_features(img, imgSize);
    
    % Classify

    label = predict(svmModel5, features);
 
end
    
\end{verbatim}

\begin{verbatim}


function features = extract_character_features(img, imageSize)
    
    if size(img,3) == 3
                img = rgb2gray(img);
    end
    img = imresize(img, imageSize);
    img = im2single(img);

    bin = imbinarize(img);
    se = strel('disk', 1);
    bin = imclose(bin, se);
    bin = bwareaopen(bin, 5);
    
    % 1. HOG 
    hog_feature = extractHOGFeatures(img, 'CellSize', [4 4]);
    
    % 2. Density global
    density = sum(bin(:)) / numel(bin);
    
    % 3. Número de agujeros
    cc = bwconncomp(~bin);
    numHoles = cc.NumObjects;
    
    % 4. Perfiles de proyección completos (16+16)
    targetLen = 16;
    hProfile = sum(bin, 1);
    vProfile = sum(bin, 2);
    hProf_norm = imresize(hProfile / (max(hProfile) + eps), [1, targetLen]);
    vProf_norm = imresize(vProfile'/ (max(vProfile)  + eps), [1, targetLen]);
    projFeatures = [hProf_norm, vProf_norm];
    
    % 5. Zoning 4x4 (densidad por zona)
    zoningF = zoningFeatures(bin, 4, 4);

    % 6. Zoning asimétrico 6x3 
    zoningF2 = zoningFeatures(bin, 6, 3);

    % 7. Crossing numbers por fila
    crossings = zeros(1, targetLen);
    binResized = imresize(bin, [targetLen, targetLen]) > 0.5;
    for row = 1:targetLen
        line = binResized(row, :);
        crossings(row) = sum(diff(line) == 1);
    end
    crossings = crossings / (max(crossings) + eps);
    
    features = [
        hog_feature, ...
        density, ...
        numHoles, ...
        projFeatures, ...
        zoningF, ...
        zoningF2, ...
        crossings
        ];
end


function zf = zoningFeatures(bin, rows, cols)
    [H, W] = size(bin);
    zf = zeros(1, rows * cols);
    idx = 1;
    for r = 1:rows
        for c = 1:cols
            rStart = round((r-1)*H/rows) + 1;
            rEnd   = round(r*H/rows);
            cStart = round((c-1)*W/cols) + 1;
            cEnd   = round(c*W/cols);
            zone = bin(rStart:rEnd, cStart:cEnd);
            zf(idx) = sum(zone(:)) / numel(zone);
            idx = idx + 1;
        end
    end
end

















%function features = extract_character_features(img, imageSize)

    % Preprocess
 
    %if size(img,3) == 3
     %   img = rgb2gray(img);
    %end
    
   % img = imresize(img, imageSize);
   % img = im2single(img);
    
  
    % BINARIZATION (para shape features)
   % bin = imbinarize(img);
    
  
    % HOG FEATURES (base)
  %  hog_feature = extractHOGFeatures(img, 'CellSize', [8 8]);
    
 
    % 1. SHAPE FEATURES (muy importantes para 8 vs B, 1 vs I)
    %density = sum(bin(:)) / numel(bin);
    
    % agujeros (topología básica)
   % cc = bwconncomp(~bin);
   % numHoles = cc.NumObjects;
    
    % =========================================================
    % 2. PROJECTION FEATURES (OCR clásico potente)
    % =========================================================
    
   % hProfile = sum(bin, 1);
   % vProfile = sum(bin, 2);
    
   % projFeatures = [
   %     mean(hProfile), std(hProfile), ...
   %     mean(vProfile), std(vProfile)
    %    ];
    
    % =========================================================
    % 3. GRADIENT FEATURES (complemento a HOG)
    % =========================================================
    
   % [Gx, Gy] = imgradientxy(img);
   % [Gmag, Gdir] = imgradient(Gx, Gy);
    
   % gradFeatures = [
     %   mean(Gmag(:)), ...
     %   std(Gmag(:))
     %   ];
    
   % dirHist = histcounts(Gdir(:), 8);
   % dirHist = dirHist / (sum(dirHist) + eps);
    
    % =========================================================
    % FINAL FEATURE VECTOR
    % =========================================================
   % features = [
  %      hog_feature, ...
  %      density, ...
  %      numHoles, ...
  %      projFeatures, ...
   %     gradFeatures, ...
  %      dirHist
   %     ];

%end

%function features = extract_character_features(img, imageSize)

    % Grayscale
    
    %if size(img,3) == 3
        %img = rgb2gray(img);
    %end
    
    % Aspect and Area Ratio
    
    %[h, w] = size(img);
    %aspectRatio = w / h;
    %areaRatio = (h * w);
    
    % Connected components
    
    %bw = imbinarize(img);
    %cc = bwconncomp(bw);
    %numComponents = cc.NumObjects;
    
    % Resize
    
    %img = imresize(img, imageSize);
    %img = im2uint8(im2double(img));
    
    % HOG FEATURES
    
    %hog_feature = extractHOGFeatures(img,'CellSize', [8 8]);
    
    %features = [hog_feature];

%end
    
\end{verbatim}

\begin{verbatim}

    
function [plateText, charLabels] = recognize_plate(aligned, blobs)
    
    plateText = "";
    charLabels = strings(size(blobs,1),1);
    
    for k = 1:size(blobs,1)
    
        bb = blobs(k,:);
    
        charImg = imcrop(aligned, bb);
    
        predictedChar = classify_character(charImg);
        
        charLabels(k) = string(predictedChar);
    
        plateText = plateText + charLabels(k);
    
    end
end
    
\end{verbatim}

\begin{verbatim}

clear;
clc;
close all;

%% Configuration

datasetPath = fullfile(pwd, "Datasets/charactersDataset2/training_data");

imds = imageDatastore(datasetPath, ...
    "IncludeSubfolders", true, ...
    "LabelSource", "foldernames");

[imdsTrain, imdsTest] = splitEachLabel(imds, 0.8, "randomized");


%% Features Extraction and Model Train

numTrain = numel(imdsTrain.Files);
numTest  = numel(imdsTest.Files);

XTrain = [];
YTrain = imdsTrain.Labels;

XTest  = [];
YTest  = imdsTest.Labels;

imgSize = [32 32];


% TRAIN FEATURES

for i = 1:numTrain

    img = readimage(imdsTrain, i);

    feat = extract_character_features(img, imgSize);

    XTrain = [XTrain; feat];

end


% TEST FEATURES 

for i = 1:numTest

    img = readimage(imdsTest, i);

    feat = extract_character_features(img, imgSize);

    XTest = [XTest; feat];

end


%% Train model 

svmModel5 = fitcecoc(XTrain, YTrain);


%% Model Evaluation

YPred = predict(svmModel5, XTest);

accuracy = sum(YPred == YTest) / numel(YTest);

fprintf("Accuracy: %.2f%%\n", accuracy * 100);

figure;
confusionchart(YTest, YPred);
title("Confusion Matrix - OCR SVM");

save("ocr_model5.mat", "svmModel5", "imgSize");
    
\end{verbatim}
